%% ========================================================================
%% Chapter 2: Manajemen Perangkat Keras & Perintah Dasar
%% ========================================================================

\chapter{Manajemen Perangkat Keras \& Perintah Dasar Sistem Operasi}
\label{ch:hardware-and-basic-commands}

\begin{abstract}
Bab ini membahas cara Linux mendeteksi dan mengelola perangkat keras, modul kernel,
driver perangkat, serta perintah-perintah dasar terminal Linux yang esensial.
\end{abstract}

%% ------------------------------------------------------------------------
\section{Deteksi Perangkat Keras di Linux}
\label{sec:hardware-detection}

% [Konten akan diisi di sini]

%% ------------------------------------------------------------------------
\section{Modul Kernel dan Driver Perangkat}
\label{sec:kernel-modules}

% [Konten akan diisi di sini]

%% ------------------------------------------------------------------------
\section{Sistem File dan /dev di Linux}
\label{sec:dev-filesystem}

% [Konten akan diisi di sini]

%% ------------------------------------------------------------------------
\section{Perintah Dasar Terminal Linux}
\label{sec:basic-commands}

% [Konten akan diisi di sini]

%% ------------------------------------------------------------------------
\section{Manipulasi Teks}
\label{sec:text-manipulation}

\subsection{Perintah sed}
\subsection{Perintah awk}
\subsection{Perintah grep}

% [Konten akan diisi di sini]

%% ------------------------------------------------------------------------
\section{Manajemen Proses}
\label{sec:process-commands}

\subsection{Perintah ps}
\subsection{Perintah top}
\subsection{Perintah kill}

% [Konten akan diisi di sini]

%% ------------------------------------------------------------------------
\section{Pemantauan Sistem}
\label{sec:system-monitoring}

% [Konten akan diisi di sini]

%% ------------------------------------------------------------------------
\section{Rangkuman}
\section{Latihan}
\section{Referensi Tambahan}
